\chapter{Atrial Fibrillation}
\index{Atrial Fibrillation}

\section*{Definition and Overview}
Atrial fibrillation (AF) is the most common serious heart rhythm disorder. In a normal heart, the upper chambers (atria) contract in a coordinated way before the lower chambers (ventricles) pump blood. In AF, the electrical signals in the atria become chaotic, so the atria quiver instead of contracting properly and the heartbeat becomes fast and irregular. AF is not immediately life-threatening in itself, but it increases the risk of stroke and can cause symptoms such as palpitations, breathlessness and fatigue. It is a long-term condition, but it can usually be well controlled with medicines and, in some people, with procedures such as catheter ablation.

\section*{Epidemiology}
Atrial fibrillation is very common and becomes more so with age; roughly one in ten people over 80 has it. It is slightly more common in men, and the number of people affected is rising as the population ages and as risk factors such as obesity and high blood pressure increase. AF is more common in people with heart disease, high blood pressure, diabetes, obstructive sleep apnoea, an overactive thyroid or heavy alcohol use. About one in three people with AF has no symptoms, and the condition is found by chance, for example during a routine health check or a screening pulse check.

\section*{Aetiology and Pathophysiology}
In AF, the normal pacemaker of the heart is bypassed by rapid, disorganised electrical signals, often originating in the pulmonary veins that enter the left atrium.

\begin{itemize}
    \item \textbf{Chaotic electrical activity:} many sites in the atria fire at once, producing rapid, irregular impulses
    \item \textbf{Irregular pumping:} the ventricles beat irregularly and sometimes too fast, so the heart pumps less efficiently
    \item \textbf{Structural changes:} stretching, scarring or enlargement of the atria from high blood pressure, valve disease or heart failure makes AF more likely
    \item \textbf{Triggering factors:} alcohol, stimulants, stress, infection, surgery and an overactive thyroid can provoke episodes
    \item \textbf{Genetics:} some families carry genes that increase susceptibility
    \item \textbf{Stroke risk:} blood pools in the quivering atrium and can form a clot, which may travel to the brain
    \item \textbf{Patterns:} AF may be paroxysmal (comes and goes), persistent or permanent
\end{itemize}

\section*{Clinical Features}
Some people feel every heartbeat, while others are unaware they have AF until a complication occurs.

\begin{itemize}
    \item Palpitations: a sensation of the heart racing, fluttering or being irregular
    \item Shortness of breath, especially with activity
    \item Fatigue and reduced exercise tolerance
    \item Dizziness or light-headedness, and sometimes fainting
    \item Chest discomfort or a feeling of pressure, not always pain
    \item Reduced ability to do everyday tasks
    \item No symptoms at all, in about a third of people
\end{itemize}

\section*{Differential Diagnosis}
An irregular or rapid heartbeat can have other causes that need to be distinguished from AF.

\begin{itemize}
    \item \textbf{Supraventricular tachycardias:} other fast, regular rhythms arising above the ventricles
    \item \textbf{Atrial flutter:} a fast but more regular rhythm that often coexists with AF
    \item \textbf{Ventricular arrhythmias:} serious fast rhythms arising in the lower chambers
    \item \textbf{Ectopic beats:} extra heartbeats (extrasystoles) that cause missed or fluttering beats
    \item \textbf{Thyroid overactivity:} a rapid, regular or irregular heart rate with tremor and weight loss
    \item \textbf{Anxiety and panic attacks:} palpitations accompanied by psychological symptoms
    \item \textbf{Stimulants:} caffeine, alcohol or some medications causing a fast heart rate
\end{itemize}

\section*{When to Seek Medical Care}
See a doctor promptly if the heart feels fast or irregular, especially if this is new, if it comes with breathlessness, dizziness or chest discomfort, or if it keeps recurring.

\textbf{Contact your local emergency services for:}
\begin{itemize}
    \item Chest pain, pressure or tightness that does not settle
    \item Severe breathlessness or difficulty breathing
    \item Fainting or near-fainting
    \item Sudden weakness, numbness, confusion or difficulty speaking (possible stroke)
    \item A very fast heart rate that will not settle
\end{itemize}

\section*{Diagnosis and Investigations}
AF is confirmed by recording the heart rhythm, and further tests look for underlying causes and assess stroke risk.

\begin{itemize}
    \item \textbf{Electrocardiogram (ECG):} the tracing shows the characteristic irregular rhythm and confirms AF
    \item \textbf{Holter monitor:} a portable recorder capturing the rhythm over 24 hours or longer, useful for episodes that come and go
    \item \textbf{Event monitors and smart devices:} can record the rhythm during symptoms
    \item \textbf{Echocardiogram:} ultrasound of the heart to check valve function and pumping strength
    \item \textbf{Blood tests:} thyroid function, kidney function, electrolytes and blood counts
    \item \textbf{Chest X-ray:} to look for heart enlargement or fluid in the lungs
    \item \textbf{Stroke risk assessment:} calculation of personal risk based on age, sex and coexisting conditions
\end{itemize}

\section*{Management}
Management has two main goals: controlling the heart rate or restoring a normal rhythm, and preventing stroke. The plan depends on symptoms, the duration of AF and individual risk factors.

\begin{itemize}
    \item \textbf{Rate control:} medicines such as beta-blockers or calcium channel blockers slow the heart rate and reduce symptoms
    \item \textbf{Rhythm control:} medicines to restore and maintain a normal rhythm, or cardioversion (a controlled electric shock)
    \item \textbf{Catheter ablation:} a procedure that destroys the tiny areas of heart muscle producing the abnormal signals
    \item \textbf{Anticoagulant medicines:} blood thinners that prevent clot formation and reduce stroke risk in most people with AF
    \item \textbf{Treat the triggers:} control blood pressure, thyroid disease, sleep apnoea, obesity and alcohol use
    \item \textbf{Lifestyle measures:} regular exercise, a heart-healthy diet and limiting alcohol
    \item \textbf{Regular review:} monitoring of symptoms, heart rate and medicines, with dose adjustment as needed
\end{itemize}

\section*{Complications}
The main danger of AF is not the irregular rhythm itself but the complications it allows.

\begin{itemize}
    \item \textbf{Stroke:} clots forming in the left atrium can travel to the brain; AF roughly triples the risk of stroke
    \item \textbf{Heart failure:} long-term rapid or inefficient pumping can weaken the heart
    \item \textbf{Reduced quality of life:} fatigue, breathlessness and exercise limitation
    \item \textbf{Dementia and cognitive decline:} possibly related to small, silent strokes and reduced blood flow
    \item \textbf{Rapid worsening of symptoms:} if the heart rate stays very fast for long periods
    \item \textbf{Treatment complications:} bleeding with blood thinners, or side effects of rhythm-control medicines and procedures
\end{itemize}

\section*{Prognosis and Outcomes}
Atrial fibrillation is a chronic condition, but most people manage it well and continue an active life. Symptoms often improve substantially with rate or rhythm control, and stroke risk is greatly reduced by anticoagulation in people who are eligible. The rhythm may return despite treatment, and some people need more than one procedure or a change of medicines over time. The main factors influencing long-term outlook are age, the presence of other heart disease, and how well stroke prevention is managed. Overall, with good care, the extra risk from AF can be largely controlled.

\section*{Prevention and Self-Care}
The underlying conditions that lead to AF can often be prevented or controlled, even though AF itself cannot always be avoided.

\begin{itemize}
    \item Control blood pressure and cholesterol; treat diabetes and sleep apnoea
    \item Limit alcohol, and avoid binge drinking, which can trigger episodes
    \item Keep a healthy weight and stay physically active
    \item Eat a heart-healthy diet and avoid smoking
    \item Manage stress and reduce stimulants such as excess caffeine
    \item Take anticoagulants exactly as prescribed, and report any unusual bleeding
    \item Seek medical attention promptly for new or worsening palpitations, breathlessness or fainting
    \item Attend scheduled heart rhythm reviews and blood tests
\end{itemize}

\section*{Sources and Further Reading}
The following reputable organisations provide reliable, current information on atrial fibrillation and stroke prevention.

\begin{itemize}
    \item National Health Service. \textit{Information on atrial fibrillation, symptoms and treatment.} https://www.nhs.uk/conditions/atrial-fibrillation/
    \item National Institute for Health and Care Excellence. \textit{Clinical guidance on the management of atrial fibrillation.} https://www.nice.org.uk/guidance/ng196
    \item Mayo Clinic. \textit{Overview of atrial fibrillation, causes and treatments.} https://www.mayoclinic.org/diseases-conditions/atrial-fibrillation
    \item World Health Organization. \textit{Fact sheets on cardiovascular diseases and stroke prevention.} https://www.who.int/health-topics/cardiovascular-diseases
    \item MedlinePlus. \textit{Health information on atrial fibrillation.} https://medlineplus.gov/atrialfibrillation.html
    \item Heart Foundation Australia. \textit{Information on arrhythmias and heart health.} https://www.heartfoundation.org.au
\end{itemize}
